\documentclass[12pt,a4paper]{article}

\usepackage[spanish]{babel}
\usepackage[utf8]{inputenc}
\usepackage[T1]{fontenc}
\usepackage{geometry}
\usepackage{xcolor}
\usepackage{hyperref}
\usepackage{graphicx}
\usepackage{float}
\usepackage{array}
\usepackage{longtable}
\usepackage{booktabs}
\usepackage{enumitem}
\usepackage{fancyhdr}
\usepackage{listings}
\usepackage{caption}

\geometry{margin=2.5cm}
\setlength{\parskip}{0.65em}
\setlength{\parindent}{0pt}

\definecolor{EduGreen}{RGB}{18,126,107}
\definecolor{EduLight}{RGB}{238,247,244}
\definecolor{EduGray}{RGB}{245,247,247}

\hypersetup{
    colorlinks=true,
    linkcolor=EduGreen,
    urlcolor=EduGreen,
    citecolor=EduGreen,
    pdftitle={Manual de Usuario EduPKIManager},
    pdfauthor={EduPKIManager}
}

\pagestyle{fancy}
\fancyhf{}
\lhead{EduPKIManager}
\rhead{Manual de Usuario}
\cfoot{\thepage}

\lstdefinestyle{terminal}{
    basicstyle=\ttfamily\small,
    breaklines=true,
    frame=single,
    backgroundcolor=\color{EduGray},
    columns=fullflexible
}

\newcommand{\nota}[1]{
    \begin{center}
    \fcolorbox{EduGreen}{EduLight}{
        \begin{minipage}{0.9\textwidth}
        \textbf{Nota.} #1
        \end{minipage}
    }
    \end{center}
}

\newcommand{\captura}[3]{
    \begin{figure}[H]
        \centering
        \IfFileExists{img_latex/#1.png}{
            \includegraphics[width=0.92\textwidth,height=0.62\textheight,keepaspectratio]{img_latex/#1.png}
        }{
            \fbox{
                \begin{minipage}[c][5cm][c]{0.88\textwidth}
                    \centering
                    \textbf{#2}\\[0.4cm]
                    \small Falta la imagen: \texttt{img\_latex/#1.png}
                \end{minipage}
            }
        }
        \caption{#3}
    \end{figure}
}

\begin{document}

\begin{titlepage}
    \centering
    \vspace*{1.7cm}
    {\Huge\bfseries Universidad La Salle - Arequipa\par}
    \vspace{1cm}
    {\Large Facultad de Ingenieria\par}
    \vspace{0.5cm}
    {\Large Escuela Profesional de Ingenieria de Software\par}

    \vspace{3.0cm}
    {\Huge\bfseries Manual de Usuario\par}
    \vspace{0.6cm}
    {\huge\bfseries EduPKIManager\par}
    \vspace{0.4cm}
    {\Large Plataforma PKI educativa\par}

    \vspace{1.2cm}
    {\Large\textbf{Proyecto 9:} PKI, Certificados X.509, Firmas Digitales y Criptografia Asimetrica\par}

    \vfill
    {\Large\textbf{Curso:} Ciberseguridad\par}
    \vspace{0.35cm}
    {\Large\textbf{Semestre:} X\par}
    \vspace{0.35cm}
    {\Large\textbf{Docente:} Marco Antonio Camacho Alatrista\par}
    \vspace{0.35cm}
    {\Large\textbf{Estudiante:} [Nombre del estudiante]\par}

    \vspace{1.1cm}
    {\Large Arequipa, Peru\par}
    {\Large 12 de julio de 2026\par}
\end{titlepage}

\tableofcontents
\newpage

\section{Introduccion}

\subsection{Proposito del Manual}

Este manual proporciona una guia paso a paso para utilizar EduPKIManager desde la perspectiva de un administrador y de un usuario final. El documento explica el acceso al sistema, la navegacion principal, la gestion de certificados digitales, la firma de documentos PDF, la verificacion criptografica, la consulta de CRL/OCSP, la prueba TLS 1.3 y la revision de auditoria.

\subsection{Descripcion General del Sistema}

EduPKIManager es una plataforma web de Infraestructura de Clave Publica orientada a entornos educativos y empresariales. El sistema implementa una Autoridad Certificadora propia, emite certificados X.509 v3, gestiona su ciclo de vida, publica listas de revocacion, responde consultas OCSP, firma documentos PDF y permite validar la trazabilidad completa de la cadena de confianza.

El sistema cuenta con dos experiencias de uso:

\begin{itemize}[leftmargin=*]
    \item \textbf{Administrador:} gestiona todos los certificados, opera la CA, revisa auditoria y ejecuta pruebas TLS.
    \item \textbf{Usuario final:} solicita certificados personales, firma documentos y valida certificados o firmas propias.
\end{itemize}

\section{Acceso al Sistema}

\subsection{Requisitos Previos}

Para usar EduPKIManager localmente se requiere:

\begin{itemize}[leftmargin=*]
    \item Docker Desktop instalado y en ejecucion.
    \item Navegador moderno: Google Chrome, Microsoft Edge, Firefox o similar.
    \item Proyecto descargado o clonado en el equipo.
    \item Archivo \texttt{hosts} configurado si se desea usar el dominio local \url{https://edupkimanager.com}.
\end{itemize}

\subsection{Acceso mediante dominio local HTTPS}

El sistema esta preparado para abrirse mediante el dominio local \url{https://edupkimanager.com}. Este dominio no necesita estar comprado, porque se configura en el archivo \texttt{hosts} del equipo para apuntar a \texttt{127.0.0.1}.

\begin{enumerate}[leftmargin=*]
    \item Abra PowerShell como administrador.
    \item Ingrese a la carpeta del proyecto.
    \item Ejecute el helper de configuracion local:
\end{enumerate}

\begin{lstlisting}[style=terminal]
.\scripts\setup_windows_tls.ps1
\end{lstlisting}

\begin{enumerate}[leftmargin=*,resume]
    \item Para que el navegador confie en el certificado HTTPS emitido por la CA propia, ejecute:
\end{enumerate}

\begin{lstlisting}[style=terminal]
.\scripts\setup_windows_tls.ps1 -TrustRootCa
\end{lstlisting}

\nota{Si PowerShell bloquea la ejecucion de scripts, ejecute primero \texttt{Set-ExecutionPolicy -Scope CurrentUser RemoteSigned} y vuelva a correr el helper.}

\subsection{Acceso local con Docker Compose}

Desde la raiz del proyecto, levante todos los servicios con:

\begin{lstlisting}[style=terminal]
docker compose up -d --build
\end{lstlisting}

Luego abra el navegador en una de las siguientes direcciones:

\begin{itemize}[leftmargin=*]
    \item Panel web HTTPS: \url{https://edupkimanager.com}
    \item Panel web HTTP de respaldo: \url{http://edupkimanager.com}
    \item Panel web por puerto local: \url{http://localhost:3000}
    \item API directa: \url{http://127.0.0.1:8000/api}
\end{itemize}

\captura{pantalla_inicio_edupkimanager}{Pantalla de inicio de EduPKIManager}{Acceso al sistema EduPKIManager}

\subsection{Credenciales de Demostracion}

\begin{longtable}{>{\bfseries}p{4cm}p{4cm}p{6cm}}
\toprule
Rol & Credencial & Uso principal \\
\midrule
Administrador & \texttt{admin / admin123} & Gestion completa de certificados, auditoria, TLS y operaciones administrativas. \\
Usuario final & \texttt{user / user123} & Solicitud de certificados personales, firma PDF y validacion de confianza. \\
\bottomrule
\end{longtable}

\nota{En un despliegue productivo se deben cambiar las credenciales demo desde el archivo \texttt{.env}.}

\section{Navegacion del Sistema}

\subsection{Inicio de Sesion}

Al abrir el sistema se presenta una pantalla de autenticacion. El usuario debe ingresar su nombre de usuario y contrasena. Despues de iniciar sesion, el frontend recibe un token y lo usa para comunicarse con la API.

\captura{formulario_inicio_sesion}{Formulario de inicio de sesion}{Inicio de sesion con rol administrador o usuario}

\subsection{Barra Superior}

La barra superior muestra el nombre EduPKIManager, el subtitulo de la consola y las opciones disponibles segun el rol autenticado.

\begin{itemize}[leftmargin=*]
    \item Si el rol es \textbf{admin}, aparecen las vistas \textit{Administracion} y \textit{Operaciones PKI}.
    \item Si el rol es \textbf{user}, aparecen las vistas \textit{Mi portal} y \textit{Firmas y validacion}.
    \item El boton de salida permite cerrar la sesion actual y volver al login.
\end{itemize}

\captura{barra_superior_roles}{Barra superior por roles}{Navegacion principal del sistema}

\subsection{Tarjetas de Resumen}

En la parte superior del panel se muestran indicadores rapidos:

\begin{itemize}[leftmargin=*]
    \item Total de certificados visibles para el rol actual.
    \item Certificados activos.
    \item Certificados revocados o suspendidos.
\end{itemize}

Para el administrador, los indicadores se calculan sobre todos los certificados. Para el usuario final, se calculan solo sobre sus certificados propios.

\section{Panel de Administrador}

\subsection{Descripcion}

El panel de administrador permite administrar el ciclo de vida completo de certificados emitidos por EduPKIManager. Desde esta vista se pueden emitir certificados para usuarios, servidores o dispositivos, consultar su estado, renovarlos, suspenderlos o revocarlos.

\captura{panel_administracion}{Panel de Administracion}{Panel de gestion completa para administradores}

\subsection{Emitir un Certificado}

Para emitir un certificado:

\begin{enumerate}[leftmargin=*]
    \item Inicie sesion con el usuario \texttt{admin}.
    \item Ingrese a la vista \textit{Administracion}.
    \item Complete el nombre comun o \textit{Common Name}; por ejemplo, \texttt{portal.edu.local}.
    \item Seleccione el tipo de certificado: usuario, servidor o dispositivo.
    \item Indique el propietario del certificado.
    \item Defina los SAN si corresponde, separados por coma.
    \item Seleccione el algoritmo de clave: RSA-2048, RSA-4096 o ECDSA P-256.
    \item Presione el boton de emision.
\end{enumerate}

Al finalizar, el sistema registra el certificado, muestra sus datos y conserva la trazabilidad de auditoria.

\captura{formulario_emision_certificados}{Formulario de emision de certificados}{Emision de certificado X.509 v3}

\subsection{Listar y Consultar Certificados}

La tabla de certificados permite revisar:

\begin{itemize}[leftmargin=*]
    \item Identificador interno.
    \item Numero de serie.
    \item Common Name.
    \item Tipo de certificado.
    \item Propietario.
    \item Estado: emitido, suspendido, revocado o renovado.
    \item Fechas de vigencia.
\end{itemize}

\subsection{Renovar un Certificado}

La renovacion marca el certificado original como renovado y crea un nuevo certificado con el mismo propietario y tipo.

\begin{enumerate}[leftmargin=*]
    \item Ubique el certificado en la tabla.
    \item Presione la accion \textit{Renovar}.
    \item Confirme la operacion si el navegador lo solicita.
    \item Verifique que aparezca un nuevo certificado activo.
\end{enumerate}

\subsection{Suspender un Certificado}

La suspension deja el certificado en estado \texttt{suspended} y lo incorpora a la CRL con la razon \texttt{certificate\_hold}.

\begin{enumerate}[leftmargin=*]
    \item Ubique el certificado que desea suspender.
    \item Presione la accion \textit{Suspender}.
    \item Revise que el estado cambie a \texttt{suspended}.
    \item Descargue la CRL para comprobar que el certificado aparece como retenido.
\end{enumerate}

\subsection{Revocar un Certificado}

La revocacion invalida el certificado y lo incluye en la CRL. Debe usarse cuando la clave privada fue expuesta, el propietario ya no debe usar el certificado o el certificado fue emitido por error.

\begin{enumerate}[leftmargin=*]
    \item Ubique el certificado.
    \item Presione \textit{Revocar}.
    \item Indique o acepte la razon de revocacion.
    \item Confirme que el estado cambie a \texttt{revoked}.
\end{enumerate}

\section{Portal de Usuario Final}

\subsection{Descripcion}

El portal de usuario final esta orientado a operaciones personales. El usuario puede solicitar certificados propios, descargarlos, consultar su estado, firmar documentos y validar certificados o firmas.

\captura{portal_usuario_final}{Portal de usuario final}{Vista Mi portal para usuario final}

\subsection{Solicitar Certificado Personal}

Para solicitar un certificado:

\begin{enumerate}[leftmargin=*]
    \item Inicie sesion con el usuario \texttt{user}.
    \item Ingrese a \textit{Mi portal}.
    \item Escriba el Common Name del certificado, por ejemplo \texttt{student.edu.local}.
    \item Presione el boton de solicitud.
    \item Verifique que el certificado aparezca en la lista de certificados propios.
\end{enumerate}

\nota{El usuario final solo puede emitir certificados de tipo usuario y solo puede operar sobre certificados cuyo propietario coincida con su cuenta.}

\subsection{Descargar Certificado Propio}

El usuario puede descargar el certificado PEM para usarlo en validaciones o pruebas externas. La clave privada solo se entrega al momento de la emision y debe custodiarse con cuidado.

\subsection{Consultar Estado}

El estado del certificado permite saber si esta activo, suspendido, revocado o renovado. Antes de firmar documentos se recomienda confirmar que el certificado este en estado \texttt{issued}.

\section{Operaciones PKI}

\subsection{Firma Digital de PDF}

EduPKIManager permite firmar documentos PDF mediante certificados emitidos por la CA propia. El sistema soporta dos formatos:

\begin{itemize}[leftmargin=*]
    \item \textbf{Firma JSON desprendida:} devuelve un sobre JSON con hash, certificado, cadena y firma.
    \item \textbf{Firma PAdES embebida:} devuelve un PDF firmado descargable.
\end{itemize}

Proceso general:

\begin{enumerate}[leftmargin=*]
    \item Ingrese a \textit{Operaciones PKI} o \textit{Firmas y validacion}.
    \item Seleccione un certificado activo.
    \item Cargue el archivo PDF.
    \item Presione \textit{Firmar JSON} o \textit{Firmar PAdES}.
    \item Descargue el resultado generado.
\end{enumerate}

\captura{modulo_firma_pdf}{Modulo de firma PDF}{Firma digital de documentos PDF}

\subsection{Verificacion de PDF}

Para verificar una firma JSON desprendida:

\begin{enumerate}[leftmargin=*]
    \item Cargue el PDF original.
    \item Cargue el archivo JSON de firma.
    \item Presione \textit{Verificar JSON}.
    \item Revise el resultado de integridad, firma, cadena de confianza, CRL y OCSP.
\end{enumerate}

Para verificar un PDF PAdES:

\begin{enumerate}[leftmargin=*]
    \item Cargue el PDF firmado.
    \item Presione \textit{Verificar PAdES}.
    \item Revise si la firma esta intacta, confiable y asociada a la Root CA de EduPKIManager.
\end{enumerate}

\captura{modulo_verificacion_pdf}{Modulo de verificacion PDF}{Verificacion de firmas desprendidas y PAdES}

\subsection{Validacion de Confianza X.509}

El modulo de confianza X.509 revisa la cadena completa del certificado:

\begin{itemize}[leftmargin=*]
    \item Firma del certificado por la Intermediate CA.
    \item Ancla de confianza Root CA.
    \item Vigencia temporal.
    \item Basic Constraints, Key Usage y Extended Key Usage.
    \item Politicas de certificado.
    \item Estado local.
    \item CRL y OCSP.
\end{itemize}

El usuario debe seleccionar el certificado y el proposito de validacion:

\begin{itemize}[leftmargin=*]
    \item Firma documental.
    \item TLS servidor.
    \item Cliente.
    \item Dispositivo.
    \item General.
\end{itemize}

\captura{modulo_confianza_x509}{Modulo de confianza X.509}{Validacion de cadena y usos permitidos}

\subsection{Consulta OCSP}

OCSP permite conocer el estado en linea de un certificado usando su numero de serie.

\begin{enumerate}[leftmargin=*]
    \item Copie el numero de serie del certificado.
    \item Pegue el valor en el modulo OCSP.
    \item Presione \textit{Consultar}.
    \item Revise el estado devuelto: \texttt{good}, \texttt{revoked} o \texttt{unknown}.
\end{enumerate}

\captura{modulo_ocsp}{Modulo OCSP}{Consulta de estado OCSP por numero de serie}

\subsection{CA y CRL}

La seccion CA y CRL permite descargar artefactos publicos:

\begin{itemize}[leftmargin=*]
    \item Certificado Root CA en PEM.
    \item CRL en PEM.
    \item CRL en DER.
    \item Manifiesto de versiones de CRL.
\end{itemize}

La CRL se actualiza automaticamente por el servicio \texttt{crl\_scheduler} y tambien puede publicarse al consultar los endpoints de CRL.

\subsection{Prueba TLS 1.3}

El sistema incluye una demostracion de handshake TLS 1.3 usando un certificado de servidor emitido por la CA propia.

\begin{enumerate}[leftmargin=*]
    \item Inicie sesion como administrador.
    \item Ingrese a \textit{Operaciones PKI}.
    \item Abra la tarjeta \textit{TLS 1.3}.
    \item Presione \textit{Probar handshake}.
    \item Confirme que cliente y servidor reporten \texttt{TLSv1.3}.
\end{enumerate}

\captura{modulo_tls13}{Modulo TLS 1.3}{Prueba de handshake TLS 1.3}

\subsection{Auditoria}

La auditoria esta disponible para el administrador. Cada operacion criptografica relevante se registra con:

\begin{itemize}[leftmargin=*]
    \item Numero de secuencia.
    \item Fecha y hora.
    \item Actor.
    \item Operacion.
    \item Resultado.
    \item Hash previo.
    \item Hash de la entrada.
\end{itemize}

El sistema verifica la cadena de hashes. Si una entrada se modifica o elimina, la auditoria deja de validar.

\captura{modulo_auditoria}{Modulo de auditoria}{Registro inmutable con hash encadenado}

\section{Flujo Recomendado de Demostracion}

Para una demostracion completa se recomienda el siguiente recorrido:

\begin{enumerate}[leftmargin=*]
    \item Iniciar sesion como administrador.
    \item Emitir un certificado de usuario.
    \item Emitir un certificado de servidor para \texttt{edupkimanager.com}.
    \item Firmar un PDF con un certificado activo.
    \item Verificar la firma del PDF.
    \item Validar la confianza X.509 del certificado.
    \item Consultar el estado OCSP.
    \item Suspender o revocar un certificado.
    \item Descargar la CRL y revisar el manifiesto.
    \item Probar el handshake TLS 1.3.
    \item Abrir auditoria y confirmar que la cadena sea integra.
\end{enumerate}

\section{Solucion de Problemas}

\subsection{El dominio \texttt{edupkimanager.com} no abre}

\begin{enumerate}[leftmargin=*]
    \item Verifique que Docker este en ejecucion.
    \item Confirme que el archivo \texttt{hosts} tenga la entrada:
\end{enumerate}

\begin{lstlisting}[style=terminal]
127.0.0.1 edupkimanager.com www.edupkimanager.com
\end{lstlisting}

\begin{enumerate}[leftmargin=*,resume]
    \item Ejecute nuevamente el helper:
\end{enumerate}

\begin{lstlisting}[style=terminal]
.\scripts\setup_windows_tls.ps1
\end{lstlisting}

\subsection{El navegador muestra advertencia de certificado}

Esto ocurre cuando la Root CA propia aun no fue importada como autoridad confiable.

\begin{lstlisting}[style=terminal]
.\scripts\setup_windows_tls.ps1 -TrustRootCa
\end{lstlisting}

Tambien puede descargar manualmente la Root CA desde:

\begin{lstlisting}[style=terminal]
http://127.0.0.1:8000/api/ca/root.pem
\end{lstlisting}

\subsection{No puedo firmar un PDF}

Revise los siguientes puntos:

\begin{itemize}[leftmargin=*]
    \item El certificado debe estar en estado \texttt{issued}.
    \item El usuario solo puede firmar con certificados propios.
    \item El archivo debe ser un PDF valido.
    \item Para PAdES, el backend debe tener disponible \texttt{pyHanko}.
\end{itemize}

\subsection{No veo todos los certificados}

El comportamiento depende del rol:

\begin{itemize}[leftmargin=*]
    \item Administrador: ve todos los certificados.
    \item Usuario final: ve solo sus certificados propios.
\end{itemize}

\subsection{La consulta OCSP devuelve \texttt{unknown}}

El estado \texttt{unknown} indica que el numero de serie consultado no corresponde a un certificado registrado en la base de datos del sistema.

\subsection{La auditoria aparece invalida}

Si la cadena de auditoria aparece invalida, alguna entrada del archivo de auditoria fue alterada, eliminada o escrita fuera del flujo normal. En ese caso se debe conservar evidencia, revisar el indice reportado y restaurar desde respaldo confiable si aplica.

\section{Buenas Practicas}

\begin{itemize}[leftmargin=*]
    \item Cambiar las credenciales demo antes de un despliegue serio.
    \item Proteger la variable \texttt{CA\_KEY\_PASSWORD}.
    \item No compartir claves privadas de certificados.
    \item Revocar certificados comprometidos inmediatamente.
    \item Revisar periodicamente la CRL y la auditoria.
    \item Usar HTTPS con TLS 1.3 para el acceso principal.
    \item Exportar evidencias de prueba antes de la entrega final.
\end{itemize}

\section{Glosario}

\begin{longtable}{>{\bfseries}p{4cm}p{10cm}}
\toprule
Termino & Definicion \\
\midrule
CA & Autoridad Certificadora encargada de emitir y firmar certificados digitales. \\
Root CA & Certificado raiz que actua como ancla de confianza. \\
Intermediate CA & CA intermedia que firma certificados finales y reduce exposicion de la Root CA. \\
Certificado X.509 & Documento digital que vincula una identidad con una clave publica. \\
CRL & Lista de Revocacion de Certificados publicada por la CA. \\
OCSP & Protocolo de consulta en linea para verificar el estado de un certificado. \\
PAdES & Perfil de firma electronica avanzada para documentos PDF. \\
TLS 1.3 & Version moderna del protocolo TLS para comunicaciones seguras. \\
SAN & Subject Alternative Name; nombres alternativos validos para un certificado. \\
\bottomrule
\end{longtable}

\end{document}
